\section{The two-sided annulus via Magnanini--Poggesi}\label{sec:trapping}

We come to the ``squeezing between two radii.'' Here the smoothness of the
inset is decisive: the inset is a \emph{smooth} domain on which the boundary
gradient $|\nabla u|$ is nearly constant
(Proposition~\ref{prop:controlled}), i.e.\ it is a near-extremal domain for
\emph{Serrin's overdetermined problem}. The quantitative stability theory
for that problem applies to smooth domains and traps the inset between two
concentric circles. We emphasise at the outset (Remark~\ref{rem:nonuniform})
that the resulting width is controlled by a constant depending on the
inset's $C^2$ geometry, which is finite (the inset is analytic) but is not,
in this soft theory, uniform in $\dT$.

We first pass from the weighted Serrin deficit \eqref{eq:serrin-inset} to
the unweighted one used in the literature.

\begin{lemma}[Unweighted Serrin deficit]\label{lem:unweighted}
Let $\Etil=E_{\hat v}$ be the inset and
$G_1:=\max_{\{u=\hat v\}}|\nabla u|<\infty$ (finite, since $\{u=\hat v\}$ is
compact in $\Om$ and $u\in C^\infty$ near it). Then, with $c=\hat\mu/\hat P$,
\begin{equation}\label{eq:unweighted}
   \eps(\Etil)^2:=\int_{\{u=\hat v\}}\bigl(|\nabla u|-c\bigr)^2\,d\HH^1
   \ \le\ G_1\,W(\Etil)\ \le\ \frac{C_0\,G_1}{4\pi}\,\dT .
\end{equation}
\end{lemma}

\begin{proof}
Pointwise on $\{u=\hat v\}$, $(|\nabla u|-c)^2=\dfrac{(|\nabla
u|-c)^2}{|\nabla u|}\,|\nabla u|\le G_1\,\dfrac{(|\nabla u|-c)^2}{|\nabla
u|}$; integrate and use \eqref{eq:serrin-inset}. Finiteness of $G_1$ is
Lemma~\ref{lem:compact-containment} (the level set lies in a fixed compact
$K_{\hat v}\Subset\Om$) together with the interior estimate
(Theorem~\ref{thm:interior-estimates}).
\end{proof}

\begin{theorem}[Two-sided annulus]\label{thm:annulus}
Let $\Etil$ be the inset and suppose, in addition to $\dT\le\delta_\star$,
that $\Etil$ is connected \textup{(}see Remark~\ref{rem:components} for the
general case\textup{)}. Then there exist a point $z\in\R^2$, radii
$0<\rho_i\le\rho_e$, an exponent $\tau\in(0,1]$, and a constant
$C_{\mathrm{MP}}$ such that
\begin{equation}\label{eq:annulus}
   B_{\rho_i}(z)\ \subset\ \Etil\ \subset\ B_{\rho_e}(z),
   \qquad
   \rho_e-\rho_i\ \le\ C_{\mathrm{MP}}\,\eps(\Etil)^{\tau}
   \ \le\ C_{\mathrm{MP}}\Bigl(\tfrac{C_0G_1}{4\pi}\Bigr)^{\tau/2}\dT^{\tau/2}.
\end{equation}
Here $\tau$ is absolute and $C_{\mathrm{MP}}$ depends only on
$\operatorname{diam}\Etil$ and on the radii of the uniform interior and
exterior ball conditions of the smooth domain $\Etil$.
\end{theorem}

\begin{proof}
The function $\hat u:=u-\hat v$ is the torsion function of the bounded
domain $\Etil$: $-\Delta\hat u=1$ in $\Etil$, $\hat u=0$ on
$\partial\Etil$, and by Lemma~\ref{lem:smooth-level} the boundary
$\partial\Etil=\{u=\hat v\}$ is a finite union of analytic Jordan curves,
so $\Etil$ is a domain of class $C^{2,\alpha}$ (indeed $C^\omega$); being
connected by hypothesis, it is a bounded $C^{2,\alpha}$ domain. On
$\partial\Etil$ the normal derivative is $\partial_\nu\hat u=-|\nabla u|$
(Lemma~\ref{lem:smooth-level}(d)), and the Serrin deficit
\[
   \Bigl(\int_{\partial\Etil}\bigl(\partial_\nu\hat u+c\bigr)^2
   d\HH^1\Bigr)^{1/2}=\eps(\Etil)
\]
is small by Lemma~\ref{lem:unweighted}, with $c=|\Etil|/P(\Etil)$ the value
forced by the trace/flux normalisation. The quantitative stability theorem
of Magnanini and Poggesi for Serrin's problem
\cite[Theorem~1.2 and \S4]{MagnaniniPoggesi2019} applies to the bounded
$C^{2,\alpha}$ domain $\Etil$: there are $z\in\Etil$,
$\rho_i=\min_{x\in\partial\Etil}|x-z|$,
$\rho_e=\max_{x\in\partial\Etil}|x-z|$, an exponent $\tau=\tau(2)\in(0,1]$,
and a constant $C_{\mathrm{MP}}$ depending only on $\operatorname{diam}
\Etil$ and the uniform interior/exterior ball radii of $\partial\Etil$,
with
\[
   \rho_e-\rho_i\ \le\ C_{\mathrm{MP}}\,\eps(\Etil)^{\tau}.
\]
The inclusion $B_{\rho_i}(z)\subset\Etil\subset B_{\rho_e}(z)$ is the
definition of $\rho_i,\rho_e$. Inserting \eqref{eq:unweighted} gives the
final bound.
\end{proof}

\begin{corollary}[The radii are pinned near $\hat r$]\label{cor:radii}
With $z,\rho_i,\rho_e$ as in Theorem~\ref{thm:annulus} and
$\hat r=\sqrt{\hat\mu/\pi}$, one has
$|\rho_i-\hat r|+|\rho_e-\hat r|\le C\bigl(\dT^{\tau/2}+\dT^{1/4}\bigr)$.
\end{corollary}

\begin{proof}
Since $B_{\rho_i}(z)\subset\Etil\subset B_{\rho_e}(z)$ and
$|\Etil|=\hat\mu=\pi\hat r^2$, we have $\pi\rho_i^2\le\pi\hat r^2\le
\pi\rho_e^2$, so $\rho_i\le\hat r\le\rho_e$; combined with $\rho_e-\rho_i\le
C\dT^{\tau/2}$ this gives $|\rho_i-\hat r|,|\rho_e-\hat r|\le
C\dT^{\tau/2}$. (The $\dT^{1/4}$ term records consistency with the
independent $L^1$ bound of Corollary~\ref{cor:L1}, which also forces the
$L^1$-optimal centre to lie within $C\dT^{1/4}$ of $z$.)
\end{proof}

\begin{remark}[The non-uniform constant: what is and is not proved]
\label{rem:nonuniform}
Theorem~\ref{thm:annulus} is the rigorous form of ``the inset is squeezed
between two radii,'' and it is genuinely two-sided. Its one limitation is
that $C_{\mathrm{MP}}$ depends on the $C^2$ geometry of $\Etil$ --- the
uniform interior/exterior ball radii of $\partial\Etil$ --- which are finite
because $\Etil$ is analytic, but which may degenerate as $\dT\to0$ if $\Om$
is rough (a deep narrow fjord of $\Om$ produces a high-curvature dent of
$\partial\Etil$, shrinking the interior ball radius and inflating
$C_{\mathrm{MP}}$). Thus the annulus \emph{width} $\to0$ for each fixed
$\Om$ as one refines, but the implied constant is not uniform over the class
$\{\dT\le\delta_\star\}$. Obtaining a \emph{uniform} annulus width
$\rho_e-\rho_i\le C\dT^{\theta}$ with absolute $C$ --- equivalently, a
Magnanini--Poggesi estimate depending only on the deficits and not on the
inset's a priori $C^2$ regularity --- is exactly the boundary-layer
uniformity that the soft methods of this paper do not supply, and is the
remaining analytic input of the larger programme. We make no claim to it
here.
\end{remark}

\begin{remark}[Connectedness and stray components]\label{rem:components}
If $\Etil$ is not connected, write $\Etil=\Etil_0\sqcup\bigsqcup_{j\ge1}C_j$
with $\Etil_0$ the component of largest area. Each $C_j$ is itself a smooth
domain and adds its full perimeter $P(C_j)\ge2\sqrt{\pi|C_j|}$ to
$P(\Etil)$; comparing with $\diso(\Etil)\le C\dT$
\eqref{eq:diso-inset} as in Lemma~\ref{lem:deficit-equiv} gives
$\sum_{j\ge1}\sqrt{|C_j|}\le C\sqrt{\hat\mu}\,\diso(\Etil)^{1/2}\le
C'\dT^{1/2}$, hence $\sum_{j\ge1}|C_j|\le C'\dT$ \textup{(}stray components
carry total area $O(\dT)$\textup{)}. Theorem~\ref{thm:annulus} then applies
to the main component $\Etil_0$, which carries area
$\hat\mu-O(\dT)$ and to which all of (i)--(v) of the main theorem transfer
verbatim; the stray components are absorbed into the $O(\dT^{1/4})$ error of
Corollary~\ref{cor:L1}. \textup{(}For $\dT\le\delta_\star$ one in fact
expects $\Etil$ connected, but we do not need this.\textup{)}
\end{remark}

\section{Synthesis}\label{sec:synthesis}

We have established every clause of Theorem~\ref{thm:main}:
\begin{itemize}[leftmargin=1.6em]
\item[(i)] analytic boundary, interior, finite perimeter, normal
$-\nabla u/|\nabla u|$: Lemmas~\ref{lem:compact-containment}
and~\ref{lem:smooth-level};
\item[(ii)] controlled isoperimetric and (weighted) Serrin deficits, both
$O(\dT)$: Propositions~\ref{prop:goodlevel}--\ref{prop:controlled} via the
identity \eqref{eq:identity} and Talenti \eqref{eq:talenti};
\item[(iii)] $L^1$-closeness to a disk, $O(\sqrt\dT)$: Corollary~\ref{cor:L1}
via FMP;
\item[(iv)] no interior holes: Proposition~\ref{prop:noholes} via
superharmonicity;
\item[(v)] inner ball / two-sided trapping between concentric circles:
Theorem~\ref{thm:annulus} via Magnanini--Poggesi, the inner radius
$\rho_i$ furnishing the inner ball.
\end{itemize}
Thus the inset of a near-extremal Saint--Venant domain is, qualitatively, a
regular (analytic) near-disk squeezed between two concentric circles, with
all of its energy- and $L^1$-scale defects controlled \emph{quantitatively
and with absolute constants} by the deficit, while the one place a
non-absolute constant enters --- the \emph{width} of the squeezing --- is
precisely the Magnanini--Poggesi constant for the smooth inset, whose
uniformisation is the sole remaining issue.

\begin{remark}[Where each hypothesis is spent]\label{rem:audit}
The interior regularity (i) uses \emph{no} smallness and \emph{no} boundary
regularity of $\Om$: it is pure interior elliptic theory. The
quantitative defect controls (ii) use only the level-set identity and
Talenti's comparison, hence hold with \emph{absolute} constants for every
admissible $\Om$. The $L^1$ closeness (iii) and the no-hole property (iv)
are likewise absolute. Only the annulus \emph{width} (v) borrows a constant
from the inset's a priori smoothness. In particular the assertion ``the
roughness of $\partial\Om$ does not reach the inset'' is exact: every
qualitative regularity feature of $\Etil$, and every quantitative feature
except the annulus width, is independent of $\partial\Om$.
\end{remark}
